% GENERATED FILE. Edit canonical lessons, then run npm run generate:books.
% canonical-source-hash: fnv1a64:b4a746bb2d97bd08
% canonical-lessons: LA-C422-divitiae, LA-C422-egeo, LA-C422-egestas, LA-C422-careo, LA-C422-abundo, LA-R422-review-421-422, LA-R422-review-final-422

\chapter{Rich and Poor}
\label{ch:la-rich-and-poor}
\hlchaptermodality{422}

\begin{chapteropening}
\textbf{By the end of this chapter:} \emph{I can talk about wealth and poverty, say that someone is in need or does without, and that a place has plenty.}

Complete the last lesson of chapter 422: I can talk about wealth and poverty, say that someone is in need or does without, and that a place has plenty.
\end{chapteropening}

\section[dīvitiae, dīvitiārum]{dīvitiae, dīvitiārum — riches, wealth}

\label{lesson:LA-C422-divitiae}



\begin{quote}
Before the new one: say the Latin for harbour dues, then the Latin for to gain.
\end{quote}

\subsection*{You'll want to know: dīvitiae, dīvitiārum}
\textbf{dīvitiae, dīvitiārum} — \textquotedblleft{}riches, wealth\textquotedblright{}.

\begin{quote}
Dīvitiae virum nōn semper fēlīcem faciunt.
\end{quote}

\begin{quote}
\emph{Riches do not always make a man happy.}
\end{quote}

\emph{dīvitiae, dīvitiārum} f. pl., built on \emph{dīves}, rich.

\begin{cousinweb}
\emph{dīvitiae} is built on \emph{dīves, dīvitis} \textquotedblleft{}rich\textquotedblright{} (EtymDB; Lewis \& Short).
\end{cousinweb}

\begin{grammarlens}[title={What you've built}]
One more word for telling, explaining and describing at B1.
\end{grammarlens}

\subsection*{Your turn}
\begin{itemize}
\raggedright
  \item \emph{Say it:} \emph{dīvitiae, dīvitiārum}
  \item \emph{Say it:} \emph{dīvitiae, dīvitiārum}, once more
  \item \emph{Say it:} a sentence of your own with \emph{dīvitiae, dīvitiārum}
  \item \emph{Recall:} say the Latin for harbour dues, then the Latin for to gain, then say \emph{dīvitiae, dīvitiārum} again
\end{itemize}

\begin{tcolorbox}[breakable,colback=teal!4,colframe=teal!35!black,title={Before you move on}]
What does dīvitiae, dīvitiārum mean? (\textquotedblleft{}Riches, wealth\textquotedblright{}.) Say it once more.
\end{tcolorbox}
\section[egeō, egēre]{egeō, egēre — to be in need (of), to lack}

\label{lesson:LA-C422-egeo}



\begin{quote}
Before the new one: say the Latin for to gain, then the Latin for riches.
\end{quote}

\subsection*{You'll want to know: egeō, egēre}
\textbf{egeō, egēre} — \textquotedblleft{}to be in need (of), to lack\textquotedblright{}.

\begin{quote}
Pecūniā egeō.
\end{quote}

\begin{quote}
\emph{I am short of money.}
\end{quote}

\emph{egeō, egēre, eguī} + ablative or genitive. \emph{Indigeō} is a stronger form; English \emph{indigent} comes from it.

\begin{cousinweb}
The origin of \emph{egeō} is not settled; Lewis \& Short give no sure account. English \emph{indigent} comes from its compound \emph{indigeō}.
\end{cousinweb}

\begin{grammarlens}[title={What you've built}]
One more word for telling, explaining and describing at B1.
\end{grammarlens}

\subsection*{Your turn}
\begin{itemize}
\raggedright
  \item \emph{Say it:} \emph{egeō, egēre}
  \item \emph{Say it:} \emph{egeō, egēre}, once more
  \item \emph{Say it:} a sentence of your own with \emph{egeō, egēre}
  \item \emph{Recall:} say the Latin for to gain, then the Latin for riches, then say \emph{egeō, egēre} again
\end{itemize}

\begin{tcolorbox}[breakable,colback=teal!4,colframe=teal!35!black,title={Before you move on}]
What does egeō, egēre mean? (\textquotedblleft{}To be in need (of), to lack\textquotedblright{}.) Say it once more.
\end{tcolorbox}
\section[egestās, egestātis]{egestās, egestātis — want, extreme poverty}

\label{lesson:LA-C422-egestas}



\begin{quote}
Before the new one: say the Latin for riches, then the Latin for to be in need.
\end{quote}

\subsection*{You'll want to know: egestās, egestātis}
\textbf{egestās, egestātis} — \textquotedblleft{}want, extreme poverty\textquotedblright{}.

\begin{quote}
In summā egestāte vīvēbat.
\end{quote}

\begin{quote}
\emph{He lived in the deepest poverty.}
\end{quote}

\emph{egestās, egestātis} f., built on \emph{egeō}, to be in need.

\begin{cousinweb}
\emph{egestās} is built on \emph{egeō} \textquotedblleft{}to be in need\textquotedblright{} (Lewis \& Short), with the ending \emph{-tās} that makes a state.
\end{cousinweb}

\begin{grammarlens}[title={What you've built}]
One more word for telling, explaining and describing at B1.
\end{grammarlens}

\subsection*{Your turn}
\begin{itemize}
\raggedright
  \item \emph{Say it:} \emph{egestās, egestātis}
  \item \emph{Say it:} \emph{egestās, egestātis}, once more
  \item \emph{Say it:} a sentence of your own with \emph{egestās, egestātis}
  \item \emph{Recall:} say the Latin for riches, then the Latin for to be in need, then say \emph{egestās, egestātis} again
\end{itemize}

\begin{tcolorbox}[breakable,colback=teal!4,colframe=teal!35!black,title={Before you move on}]
What does egestās, egestātis mean? (\textquotedblleft{}Want, extreme poverty\textquotedblright{}.) Say it once more.
\end{tcolorbox}
\section[careō, carēre]{careō, carēre — to be without, to lack}

\label{lesson:LA-C422-careo}



\begin{quote}
Before the new one: say the Latin for to be in need, then the Latin for want.
\end{quote}

\subsection*{You'll want to know: careō, carēre}
\textbf{careō, carēre} — \textquotedblleft{}to be without, to lack\textquotedblright{}.

\begin{quote}
Multī cīvēs domō carent.
\end{quote}

\begin{quote}
\emph{Many citizens have no home.}
\end{quote}

\emph{careō, carēre, caruī} + ablative: \emph{careō pecūniā}, I have no money. Unlike \emph{egeō}, it says the thing is absent without saying that you need it.

\begin{cousinweb}
\emph{careō} is inherited from Proto-Italic \emph{*kazēō} (EtymDB), from a Proto-Indo-European root meaning \textquotedblleft{}to cut\textquotedblright{}: to be cut off from a thing.
\end{cousinweb}

\begin{grammarlens}[title={What you've built}]
One more word for telling, explaining and describing at B1.
\end{grammarlens}

\subsection*{Your turn}
\begin{itemize}
\raggedright
  \item \emph{Say it:} \emph{careō, carēre}
  \item \emph{Say it:} \emph{careō, carēre}, once more
  \item \emph{Say it:} a sentence of your own with \emph{careō, carēre}
  \item \emph{Recall:} say the Latin for to be in need, then the Latin for want, then say \emph{careō, carēre} again
\end{itemize}

\begin{tcolorbox}[breakable,colback=teal!4,colframe=teal!35!black,title={Before you move on}]
What does careō, carēre mean? (\textquotedblleft{}To be without, to lack\textquotedblright{}.) Say it once more.
\end{tcolorbox}
\section[abundō, abundāre]{abundō, abundāre — to have plenty (of), to abound; to overflow}

\label{lesson:LA-C422-abundo}



\begin{quote}
Before the new one: say the Latin for want, then the Latin for to be without.
\end{quote}

\subsection*{You'll want to know: abundō, abundāre}
\textbf{abundō, abundāre} — \textquotedblleft{}to have plenty (of), to abound; to overflow\textquotedblright{}.

\begin{quote}
Urbs frūmentō abundat.
\end{quote}

\begin{quote}
\emph{The city has plenty of grain.}
\end{quote}

\emph{abundō, abundāre, abundāvī, abundātum} + ablative, built on \emph{unda}, a wave: to overflow. English \emph{abound} and \emph{abundant}.

\begin{cousinweb}
\emph{abundō} is built from \emph{ab} \textquotedblleft{}away, over\textquotedblright{} and \emph{undō} \textquotedblleft{}to surge\textquotedblright{}, from \emph{unda} \textquotedblleft{}a wave\textquotedblright{} (EtymDB; Lewis \& Short): to flow over. English \emph{abound} and \emph{abundant} come from it.
\end{cousinweb}

\begin{grammarlens}[title={What you've built}]
One more word for telling, explaining and describing at B1.
\end{grammarlens}

\subsection*{Your turn}
\begin{itemize}
\raggedright
  \item \emph{Say it:} \emph{abundō, abundāre}
  \item \emph{Say it:} \emph{abundō, abundāre}, once more
  \item \emph{Say it:} a sentence of your own with \emph{abundō, abundāre}
  \item \emph{Recall:} say the Latin for want, then the Latin for to be without, then say \emph{abundō, abundāre} again
\end{itemize}

\begin{tcolorbox}[breakable,colback=teal!4,colframe=teal!35!black,title={Before you move on}]
What does abundō, abundāre mean? (\textquotedblleft{}To have plenty (of), to abound; to overflow\textquotedblright{}.) Say it once more.
\end{tcolorbox}
\section[(dialogue)]{Laws, money and agreements}

\label{lesson:LA-R422-review-421-422}



\begin{quote}
Say the Latin for trade and to do business.
\end{quote}

\subsection*{Your turn}
Say these aloud:

\begin{itemize}
\raggedright
  \item trade, to do business, market day, harbour dues, to gain
  \item riches, to be in need, want, to be without, to have plenty
\end{itemize}

\begin{tcolorbox}[breakable,colback=teal!4,colframe=teal!35!black,title={Before you move on}]
Trade? (\textbf{commercium, commerciī}.) To do business? (\textbf{negōtior, negōtiārī}.) Market day? (\textbf{nūndinae, nūndinārum}.) Harbour dues? (\textbf{portōrium, portōriī}.) To gain? (\textbf{lucror, lucrārī}.) Riches? (\textbf{dīvitiae, dīvitiārum}.) To be in need? (\textbf{egeō, egēre}.) Want? (\textbf{egestās, egestātis}.) To be without? (\textbf{careō, carēre}.) To have plenty? (\textbf{abundō, abundāre}.)
\end{tcolorbox}
\section[(dialogue)]{Rich and Poor — once more}

\label{lesson:LA-R422-review-final-422}



\begin{quote}
Say the Latin for riches and to be in need.
\end{quote}

\subsection*{Your turn}
\begin{itemize}
\raggedright
  \item \emph{Say it:} riches, to be in need, want, to be without, to have plenty
\end{itemize}

\begin{tcolorbox}[breakable,colback=teal!4,colframe=teal!35!black,title={Before you move on}]
Riches? (\textbf{dīvitiae, dīvitiārum}.) To be in need? (\textbf{egeō, egēre}.) Want? (\textbf{egestās, egestātis}.) To be without? (\textbf{careō, carēre}.) To have plenty? (\textbf{abundō, abundāre}.)
\end{tcolorbox}
